\documentclass{article}
\usepackage[UTF8]{ctex}
\usepackage{amsmath,amssymb,amsfonts}
\usepackage{enumitem}
\usepackage{geometry}
\geometry{a4paper, margin=1in}
% === 错题本排版优化（不参与题目标识） ===
\setlength{\parindent}{0pt}
\setlength{\parskip}{0.35em plus 0.10em minus 0.08em}
\setlength{\emergencystretch}{3em}
\setlist{itemsep=0.25em, topsep=0.40em, parsep=0pt}
\allowdisplaybreaks[2]
% === 错题本排版优化结束 ===
\begin{document}
\section*{云南大学2024年数学分析真题}
\section*{193．云南大学2024年数学分析试卷}
\begin{itemize}

\subsection{题目1}
  \item[1．] 极限 \(\lim _{x \rightarrow 0}\left(\cot x-\frac{1}{x}\right)=\) \(\_\_\_\_\)．

\subsection{题目2}
  \item[2．] 设 \(f(x)=\left\{\begin{array}{ll}e^{a x}, & x \leq 0 ; \\ (x-b)^{3}, & x>0 .\end{array}\right.\) 在 \(x=0\) 处可导，则 \(a\) 和 \(b\) 的值分别是 \(\_\_\_\_\)．
\end{itemize}

\subsection{题目3}
3．不定积分 \(\int \frac{\mathrm{d} x}{2 e^{-x}+e^{x}+2}=\) \(\_\_\_\_\)．

\subsection{题目4}
4．曲面 \(z=x+2 y+\ln \left(1+x^{2}+y^{2}\right)\) 在 \((0,0,0)\) 的切平面方程是 \(\_\_\_\_\)．

\subsection{题目5}
5．设 \(f(x, y)=\int_{0}^{x y} e^{x t^{2}} \mathrm{~d} t\) ，则 \(\left.\frac{\partial^{2} f}{\partial x \partial y}\right|_{(1,1)}=\) \(\_\_\_\_\)．

\subsection{题目6}
6．曲线 \(y=\int_{0}^{x} \tan t \mathrm{~d} t, 0 \leq x \leq \frac{\pi}{4}\) 的弧长 \(s=\) \(\_\_\_\_\)．

二、（15 分）按定义证明： \(\lim _{x \rightarrow 1} \frac{x(x-1)}{x^{2}-1}=\frac{1}{2}\) ．

三．（15 分）判别级数 \(\sum_{n=1}^{\infty} \frac{a^{n} n!}{n^{n}}(a>0)\) 的敛散性。

四．（15 分）设函数 \(f(x)\) 和 \(g(x)\) 在区间 \([a, b]\) 上二阶可导，并且 \(g^{\prime \prime}(x) \neq 0, f(a)=f(b)=g(a)=\) \(g(b)=0\) ．证明：

\begin{itemize}
  \item[（1）] 在开区间 \((a, b)\) 内，\(g(x) \neq 0\) ．
  \item[（2）] 在开区间 \((a, b)\) 内，至少存在一点 \(\xi\) ，使得 \(\frac{f(\xi)}{g(\xi)}=\frac{f^{\prime \prime}(\xi)}{g^{\prime \prime}(\xi)}\) ．
\end{itemize}

五．（15 分）已知函数 \(z=u e^{a x+b y}\) ，其中 \(u=u(x, y)\) 具有二阶偏导数，且 \(\frac{\partial^{2} u}{\partial x \partial y}=0\) ．确定常数 \(a, b\) ，使得函数 \(z=z(x, y)\) 满足方程

\[
\frac{\partial^{2} z}{\partial x \partial y}+2 \frac{\partial z}{\partial x}+3 \frac{\partial z}{\partial y}+6 z=0
\]

六．（15 分）设 \(f(x)\) 在 \(x \neq 0\) 时一阶连续可导，且 \(f(1)=0\) ．求函数 \(f\left(x^{2}+y^{2}\right)\) ，使得曲线积分

\[
\int_{L} y\left(2-f\left(x^{2}+y^{2}\right)\right) \mathrm{d} x+x f\left(x^{2}+y^{2}\right) \mathrm{d} y
\]

与路径无关，其中 \(L\) 为不经过原点的光滑曲线。

七．（15 分）求级数 \(\sum_{n=1}^{\infty} n(n+2) x^{n-1}\) 的和函数．

八．（15 分）计算极限 \(\lim _{a \rightarrow+\infty} \iint_{D} \min \{x, y\} e^{-\left(x^{2}+y^{2}\right)} \mathrm{d} x \mathrm{~d} y\) ，其中 \(D:[-a, a] \times[-a, a]\) ．

九．（15 分）设 \(f(x)\) 在 \([a,+\infty)\) 上一致连续，\(g(x)\) 在 \([a,+\infty)\) 上连续，且 \(\lim _{x \rightarrow+\infty}[f(x)-g(x)]=0\) ．证明：\(g(x)\) 在 \([a,+\infty)\) 上一致连续。
\end{document}
